\documentclass[11pt]{article}
\usepackage[margin=1.2in]{geometry}
\usepackage{amsmath,amssymb,amsthm}
\usepackage{hyperref}
\hypersetup{colorlinks=true, linkcolor=blue, urlcolor=blue}
\usepackage{parskip}
\emergencystretch=3em

\newcommand{\R}{\mathbb{R}}
\newcommand{\code}[1]{\texttt{#1}}

\title{Tide retrospective: the located grand headline\\
\large grand composition, part 3}
\author{Claude (Anthropic), with GPT-5.6 Sol consults}
\date{2026-08-10}

\begin{document}
\maketitle

\begin{abstract}
The composition arc reaches the note's inverse Theorem 3.1 in located
form: two localized nondegenerate losses at \emph{unknown} centres,
one physical data premise --- superPoly agreement of the actual
translated-loss moment families over smooth compactly supported tests
in a common region containing balls around both centres --- and the
conclusion that the centres coincide and every positive-order
derivative tensor agrees
(\code{located\_positive\_jet\_recovery\_of\_ccData}). The proof is
one page: the located cutoff removal of part 2 turns the single test
premise into located-moment data for every smooth polynomial-growth
observable; the coordinate observables recover the location; and one
reparametrization identity --- once located at $p$, the located
moment of $Q(\cdot - p)$ \emph{is} the centred moment of $Q$ --- turns
the same data into centred data for every observable the
base-case-free jet headline wants. The deliberation had routed this
step through the centred compactly-supported headline, which would
have required the common region to contain that theorem's own cutoff
supports; reparametrizing into the superPoly-form headline dissolves
the second support condition entirely. Zero fix rounds.
\end{abstract}

\section{Setting}

Parts 1--2 built the translation bridge and the located cutoff
removal. What remained for a located headline was glue --- but glue
with one genuine choice in it, namely how to hand the actual-region
data to a centred recovery theorem whose own hypotheses mention
centred supports.

\section{The thing formalised}

\begin{sloppypar}
\code{located\_positive\_jet\_recovery\_of\_ccData} takes the package
families of both centred losses, the symmetry hypotheses of the
centred headline, the common-region data premise (per package order,
as in the perimeter corpus), and concludes $p_1 = p_2$ together with
positive-order jet equality. Location: part 2 at the coordinate
observables (smooth by \code{EuclideanSpace.proj}, linear growth),
then \code{location\_eq\_of\_superPoly\_first\_moments}. Jets: after
\code{subst}, for any smooth polynomial-growth $Q$ the actual
observable $Q(\cdot - p)$ is smooth and of polynomial growth
(translation preserves both), and
\code{locatedMomentT\_sub\_observable} collapses its located moment
to the centred moment of $Q$; instantiating $Q$ at coordinate
products and monomial tests feeds
\code{smooth\_positive\_jet\_recovery\_of\_superPoly\_moments}
directly.
\end{sloppypar}

\section{Where progress stalled}

Nowhere --- zero fix rounds, the second such tide in this arc. The
catalogued disciplines that used to be fix rounds are now defaults:
\code{beta\_reduce} before rewriting inside \code{SuperPoly.congr}
per-point goals, \code{sub\_eq\_neg\_add} to align translated
observables with the \code{comp\_const\_add} growth lemma, and
checking the import closure before naming sibling theorems.

\section{Mathlib lookup versus new mathematics}

Lookup: \code{EuclideanSpace.proj} smoothness, \code{subst}-driven
filter bookkeeping, nothing else new. New mathematics: none. The one
idea --- reparametrize the observable class rather than re-localize
the tests --- is elementary but load-bearing: it is what makes the
located headline need \emph{no} hypotheses beyond part 2's.

\section{What the next tide inherits}

The arc capstone (tide 4): transport the centred conclusions to the
actual losses $\Lambda_i = L_i(\cdot - p)$ --- iterated derivatives
at $p$, and for analytic losses the germ at $\mathcal{N}(p)$ with the
symmetry hypotheses discharged by
\code{AnalyticAt.iteratedFDeriv\_isSymm}. After that, the note-side
markers for the located theorems (gated, as all TeX actions are).
\end{document}
